\begin{enumerate}
	\item Migliorare il modello di simulazione
	\item Disturbi e parte predittiva
\end{enumerate}

\begin{itemize}
	\item Sviluppare algoritmo in base al giorno successivo 
	\item Diversi (modi) allocazione di risorse (???)
	\item Definire il tasso di errore (scostamento)
	\item Trovare allocazione di risorse che minimizza l'errore
	\begin{itemize}
		\item Definire l'errore
		\item Definire la funzione obiettivo
		\item Certo numero di combinazioni
		\item Simulazione "inline"
		\item Simulazione off-line sera prima come riferimento per giorno dopo
		\item orizzonte più piccolo in linea
	\end{itemize}
\end{itemize}

Si dovrebbe usare un sistema closed-loop dove il riferimento è il numero
di missioni e si dovrebbe minimizzare il tempo di esecuzione
L'errore è la differenza tra numero di missioni previste e quelle svolte\newline


\noindent BATCH FUTURO -> Sto in x\_k devo decidere x\_k+1 -> ragionare a chunk "temporali"\newline

\noindent Si usa lo stato:
\begin{itemize}
	\item missioni residue
	\item posizione delle risorse
\end{itemize}

\noindent Poi:
\begin{itemize}
	\item valutazione indice di prestazione
	\item svolgere la simulazione per una durata di tempo previsto
\end{itemize}


\section{APPUNTI 11 Settembre incontro Beta80}
\begin{enumerate}
	\item Raggruppare le missioni per numero ordine e svolgere le missione in ordine di priorità (NM\_PRIORITY)
	\item Fare eseguire le missioni relative ad un singolo ordine ad un SOLO operatore
	\item Definire una deadline per lo svolgimento di un singolo ordine in base al numero di missioni per ordine. Es. \texttt{num\_missioni*durata\_euristica}
\end{enumerate}
Raggruppare le missioni per numero d'ordine ed eseguirle in base alla priorità (NM\_PRIORITY);Far eseguire le missioni relative a un singolo ordine a un SOLO operatore;Definire una deadline per il completamento di ciascun ordine in base al numero di missioni associate, ad esempio:
numero\_missioni × durata\_media\_missione;Individuare un ordine di esecuzione delle missioni ottimale.

\subsection{Appunti prof}

\begin{enumerate}
	\item Un ordine si snoda in una o più missioni. Pertanto, a ciascuna
	missione è associato un identificativo di ordine. Tutte le missioni
	relative allo stesso ordine vanno eseguite dallo stesso operatore.
	Tuttavia, ciascuna missione ha una priorità da interpretarsi solo
	rispetto alle missioni dell'ordine di cui essa fa parte. Ad un ordine
	può essere associata una deadline (da inserire a mano, non è presente
	attualmente nei file forniti) da intendersi come durata massima di
	esecuzione dell'ordine (dall'inizio dell'esecuzione della prima missione
	al completamento dell'ultima missione presente nell'ordine).
	Si vuole ottimizzare il numero di risorse per l'esecuzioni di missioni
	in un turno di lavoro tenendo conto di questi nuovi vincoli;
	\item Rispetto all'ottimizzazione di cui al punto 1) al tempo di esecuzione
	si verificano per svariati motivi scostamenti. Si vuole rischedulare con
	una frequenza da valutare in base alla complessità computazionale
	l'ordine delle missioni calcolato offline al fine di ridurre l'errore
	rispetto ad indici tipo percentuale di utilizzo di una risorsa,
	percentuale di ordini completata sul totale etc. L'approccio da
	utilizzare è del tipo predittivo con "receding horizon". Si ipotizza di
	suddividere l'elenco di missioni in batch di dimensione fissata, ciò è
	essenziale per motivi pratici, non avrebbe senso ogni volta considerare
	l'intero residuo di di missioni da eseguire. Gli istanti sono i batch,
	il passo k è il batch appena eseguito, il passo k+1 il successivo, se
	l'orizzonte è 3, predico il comportamento e lo ottimizzo sui successivi
	3 batch e ne applico l'ordine solo al successivo e così in avanti. Il
	grado di libertà è costituito dal sequenziamento degli ordini, essendo
	stavolta fissato il numero di risorse.
\end{enumerate}
